\documentclass[twocolumn,10pt,a4paper]{article}

% Page layout
\usepackage[margin=2cm, columnsep=0.6cm]{geometry}

% Essential packages
\usepackage{amsmath,amssymb,amsthm}
\usepackage{mathtools}
\usepackage{bm}
\usepackage{graphicx}
\usepackage{hyperref}
\usepackage{natbib}
\usepackage{physics}
\usepackage{booktabs}
\usepackage{array}
\usepackage{multirow}
\usepackage{enumitem}
\usepackage{float}
\usepackage{authblk}

% Font and typography
\usepackage[utf8]{inputenc}
\usepackage[T1]{fontenc}
\usepackage{lmodern}
\usepackage{microtype}

% Theorem environments
\newtheorem{theorem}{Theorem}[section]
\newtheorem{lemma}[theorem]{Lemma}
\newtheorem{proposition}[theorem]{Proposition}
\newtheorem{corollary}[theorem]{Corollary}
\theoremstyle{definition}
\newtheorem{definition}[theorem]{Definition}
\newtheorem{axiom}{Axiom}
\newtheorem{principle}[theorem]{Principle}
\newtheorem{example}[theorem]{Example}
\theoremstyle{remark}
\newtheorem{remark}[theorem]{Remark}

% Hyperref setup
\hypersetup{
    colorlinks=true,
    linkcolor=blue,
    citecolor=blue,
    urlcolor=blue,
    pdftitle={Mass Spectrometry as Empirical Resolution of Loschmidt's Paradox},
    pdfauthor={Kundai Farai Sachikonye}
}

% Custom commands
\newcommand{\kB}{k_{\mathrm{B}}}
\newcommand{\Mcount}{M}
\newcommand{\Nstate}{N_{\mathrm{state}}}
\newcommand{\Partition}{\mathcal{P}}
\newcommand{\Operation}{\mathcal{O}}
\newcommand{\Negation}{\mathcal{N}}
\newcommand{\StateSpace}{\mathcal{S}}
\newcommand{\Real}{\mathbb{R}}
\newcommand{\Natural}{\mathbb{N}}
\newcommand{\tp}{\langle\tau_p\rangle}
\newcommand{\dM}{\mathrm{d}M}
\newcommand{\dt}{\mathrm{d}t}

\begin{document}

\title{\textbf{Mass Spectrometry as Empirical Resolution of Loschmidt's Paradox: \\
Constitutive Asymmetry, the Time-Count Identity, and the Incoherence of Time-Reversal}}

\author{
Kundai Farai Sachikonye\\
Technical University of Munich\\
Technical University of Munich Chair of Proteomics and Bioanalytics \\
Experimental Bioinformatics \\
Maximus-von-Imhof-Forum 3\\
85354 Freising, Germany\\
\texttt{kundai.sachikonye@wzw.tum.de}
}

\date{\today}

\maketitle

\begin{abstract}
\noindent Loschmidt's paradox has stood for one hundred and fifty years as a structural challenge to the foundations of statistical mechanics: if the microscopic laws of motion are time-reversible, why is the macroscopic world irreversible? Standard resolutions---Boltzmann's statistical argument, Gibbs's coarse-graining, Penrose's cosmological boundary conditions, Zurek's environmental decoherence---each address the question within a framework that presupposes time-reversibility as a coherent, falsifiable property of dynamics.

\vspace{0.2cm}
\noindent We resolve the paradox by dissolving its presupposition. Working within the partition-counting framework of bounded phase space dynamics, we establish that mass spectrometry constitutes an empirical demonstration of the impossibility---and, more sharply, the \emph{incoherence}---of time-reversal. The argument proceeds through seven principal results. (1) \emph{The Time-Count Identity:} time is the partition count $\Mcount$ generated by oscillation, related by $t = \Mcount/f$ exactly. (2) \emph{The Loop Identity Lemma:} reversibility plus determinism plus non-trivial dynamics force the observable state to carry partition structure beyond its microstate coordinates. (3) \emph{The Sliding-Endpoint Theorem:} empirical reproducibility of mass spectra is logically equivalent to the irreversibility of partition transitions. (4) \emph{The Constitutive Asymmetry Theorem:} dynamical trajectories are sequences of physical operations, each constitutively one-way; the time-reversibility of the continuous-flow description $\dot{x} = f(x,t)$ is a coarse-graining artifact. (5) \emph{The Rewind-as-Forward Principle:} any operation purporting to reverse a process is itself a forward operation that increments rather than decrements the partition count. (6) \emph{The Incoherence-of-Reversibility Theorem:} no operation can decrement the global partition count; time-reversal is therefore not a property dynamics either possesses or lacks but a category error generated by confusing automorphisms of solution sets with physical operations on systems. (7) \emph{The Source-Analyzer Indeterminacy Theorem:} the detector recovers only the fixed-point value of the analyzer chain's composed negation operator, not the operator itself.

\vspace{0.2cm}
\noindent These results converge on a single physical picture: mass, time, identity, and partition count are four coordinatisations of a single quantity, related by unit conversions through the oscillator frequency, the speed of partition propagation, and Planck's constant. A mass spectrometer does not measure a pre-existing quantity called mass; it constitutes the mass-record through a sequence of operations whose intrinsic directionality is the arrow of time itself. Every reproducible mass spectrum is a written record of constitutive asymmetry. The Second Law of Thermodynamics emerges not as a statistical tendency or an empirical regularity but as the operational signature of what time is.

\vspace{0.2cm}
\noindent We provide a comprehensive numerical validation of the principal theorems through seven experiments comprising $162$ sub-tests across approximately five orders of magnitude of physical parameters. All $162$ sub-tests pass. Identities involving fundamental constants ($t = M/f$, $E = \hbar\omega_0$) hold to floating-point precision; operational claims (irreversibility of the partition count, asymmetry of becoming, oscillation as forward re-placement) are verified across hundreds of randomised trials. The principle that ``oscillation is forward re-placement, not return'' is verified directly on a $10{,}000$-cycle simulation in which the partition count reaches $99{,}990$ monotonically---five orders of magnitude beyond what the naive-return hypothesis would predict.

\vspace{0.3cm}
\noindent\textbf{Keywords:} Loschmidt's paradox, time-reversibility, mass spectrometry, partition counting, constitutive asymmetry, arrow of time, Second Law of Thermodynamics, irreversibility
\end{abstract}

%==============================================================================
\section{Introduction}
\label{sec:introduction}
%==============================================================================

\subsection{The Paradox}

In 1876, Josef Loschmidt posed the following challenge to Ludwig Boltzmann's emerging statistical interpretation of the Second Law of Thermodynamics \citep{loschmidt1876,loschmidt1877}. The microscopic equations of motion for a gas of $N$ particles are Hamiltonian:
\begin{equation}
    \dot{q}_i = \frac{\partial H}{\partial p_i}, \qquad \dot{p}_i = -\frac{\partial H}{\partial q_i},
    \label{eq:hamilton}
\end{equation}
and these equations are invariant under the simultaneous transformation $t \to -t$, $p_i \to -p_i$. Therefore, Loschmidt argued, if a gas evolves from a low-entropy initial state at $t = 0$ to a high-entropy state at $t = T$, one could in principle reverse all molecular velocities at $t = T$ and the gas would evolve back to its low-entropy state. Entropy would decrease, in apparent violation of the Second Law.

Boltzmann's response was statistical \citep{boltzmann1877}: the reversed evolution is not impossible, merely overwhelmingly improbable. The probability of spontaneously arriving at the precise reversed state is of order $e^{-\Delta S/\kB} \sim e^{-N}$ where $N \sim 10^{23}$, an astronomically small number. The Second Law is thus a statistical regularity, not a fundamental law.

This statistical resolution has been the orthodox answer for one hundred and fifty years. It has been refined, supplemented, and reformulated by Gibbs's coarse-graining \citep{gibbs1902}, Jaynes's information-theoretic interpretation \citep{jaynes1957}, Prigogine's dissipative-structures programme \citep{prigogine1977}, Penrose's cosmological argument from low-entropy initial conditions \citep{penrose1989,penrose2004}, and Zurek's quantum decoherence approach \citep{zurek1991,zurek2003}. Each addresses a different facet of the paradox; none, we shall argue, dissolves it.

\subsection{Why Existing Resolutions Are Incomplete}

The standard resolutions all share a structural feature: they accept Loschmidt's framing of the question. They take it as given that time-reversibility of microscopic dynamics is a meaningful, falsifiable property, and they explain why---despite this property---we observe macroscopic irreversibility. The reasons offered vary:

\begin{itemize}[leftmargin=1.2em,itemsep=0.2em]
    \item \textbf{Boltzmann:} Because reversed states are improbable.
    \item \textbf{Gibbs:} Because coarse-graining loses information.
    \item \textbf{Jaynes:} Because we are uncertain about microstates.
    \item \textbf{Penrose:} Because the universe began in a low-entropy state.
    \item \textbf{Prigogine:} Because real systems are open and dissipative.
    \item \textbf{Zurek:} Because environmental entanglement destroys quantum coherence.
\end{itemize}

Each of these is correct within its proper domain. But none answers the question: \emph{what would it mean, operationally, to reverse the microscopic dynamics?} The standard answers proceed as if this question had an obvious answer---``apply $t \to -t, p \to -p$ to the state at time $T$''---and then explain why doing so is hard or improbable. We shall argue that the question itself is malformed: there is no operation in the physical world whose enactment corresponds to the mathematical operation $t \to -t$.

\subsection{The Approach: Empirical Witness via Mass Spectrometry}

We take a different approach. Rather than offering yet another within-framework resolution, we use a working laboratory technology---mass spectrometry---to expose the framework itself as inadequate. The argument is empirical in character: we show that the mere existence of reproducible mass spectra is logically incompatible with the framework in which Loschmidt's paradox arises.

The argument is structured as follows. First, we establish that mass spectrometry is fundamentally a partition-counting process: the trajectory of an ion through an analyzer chain is a sequence of countable partition transitions, and the detector reads the accumulated count \citep{sachikonye2025counting,sachikonye2025partition}. Second, we observe that mass spectrometers produce stable, reproducible peaks: the same molecular species, measured under the same conditions, yields the same $m/z$ value to within instrumental precision. Third, we show that this reproducibility, combined with the partition-counting interpretation, logically entails the irreversibility of partition transitions, the constitutive asymmetry of physical operations, and ultimately the incoherence of the very concept of time-reversal.

The conclusion is more radical than the standard resolutions: time-reversibility is not a difficult-to-achieve property of dynamics; it is not even a meaningful property. It is a category error generated by mistaking automorphisms of solution sets for physical operations on systems. Loschmidt's paradox dissolves because its premise was malformed.

\subsection{Structure of the Paper}

Section~\ref{sec:phenomenology} presents the phenomenological anchor: the bare empirical fact that mass spectrometers produce reproducible spectra, and why this fact requires explanation. Section~\ref{sec:time-count} establishes the Time-Count Identity, the most important conceptual move of the paper: time is constituted by, not embedded in, partition operations. Section~\ref{sec:loop} proves the Loop Identity Lemma: reversibility plus determinism plus non-trivial dynamics force the observable state to carry partition structure beyond microstate coordinates. Section~\ref{sec:sliding} proves the Sliding-Endpoint Theorem: the empirical reproducibility of MS is logically equivalent to partition irreversibility. Section~\ref{sec:constitutive} establishes the Constitutive Asymmetry Theorem: real trajectories are sequences of operations, each one-way, and the apparent time-reversibility of continuous-flow descriptions is a coarse-graining artifact. Section~\ref{sec:rewind} states the Rewind-as-Forward Principle and proves the Incoherence-of-Reversibility Theorem. Section~\ref{sec:negation} develops identity as the fixed point of negation, the operational content of the third route to mass. Section~\ref{sec:source} proves the Source-Analyzer Indeterminacy Theorem. Section~\ref{sec:routes} establishes the equivalence of the three routes to mass and the Mass-Time-Identity-Count Equivalence. Section~\ref{sec:resolution} states the resolution of Loschmidt's paradox. Section~\ref{sec:conclusion} concludes.

%==============================================================================
\section{Phenomenological Anchor: Why Mass Spectrometers Work}
\label{sec:phenomenology}
%==============================================================================

\subsection{The Bare Empirical Fact}

Consider the most ordinary use of a mass spectrometer. A sample is introduced, ionised, and the ions are separated by their mass-to-charge ratio $m/z$ and detected. The output is a spectrum: a set of peaks at specific $m/z$ values, each corresponding to a molecular species. If the same sample is measured again under the same conditions, the same peaks appear at the same $m/z$ values, to within the resolution of the instrument.

This is so familiar that its remarkable nature is rarely noticed. We propose that it is the most direct empirical anchor of the irreversibility of nature.

\subsection{Why Reproducibility Is Not Obvious}

A mass spectrometer is a Hamiltonian system: ions move in electromagnetic fields under the Lorentz force,
\begin{equation}
    m\ddot{\mathbf{r}} = q\left(\mathbf{E}(\mathbf{r}, t) + \dot{\mathbf{r}} \times \mathbf{B}(\mathbf{r})\right),
    \label{eq:lorentz}
\end{equation}
and the equations of motion are time-reversible: if $(\mathbf{r}(t), \dot{\mathbf{r}}(t))$ is a solution, so is $(\mathbf{r}(-t), -\dot{\mathbf{r}}(-t))$. By the Loschmidt argument, an ion arriving at the detector at time $T$ should be just as compatible with a reversed-evolution as with the forward evolution. Why, then, does the detector produce the same peak when we run the experiment again?

The standard answer (statistical mechanics) is unsatisfactory in this context: a single ion does not undergo statistical fluctuations of $10^{23}$ particles. The single-ion trajectory is deterministic and time-reversible at the level of \eqref{eq:lorentz}. Yet the spectrum is reproducible. Why?

\subsection{The Question to be Answered}

We pose the question sharply:

\begin{quote}
\emph{What feature of physical reality must hold in order for mass spectrometers to produce reproducible spectra, given that the underlying dynamics is described by time-reversible equations?}
\end{quote}

We shall show that the answer is: the partition transitions constituting the ion's journey are constitutively one-way, and the operation that purports to reverse them is itself a forward operation. The reproducibility of the spectrum is the empirical signature of this constitutive asymmetry.

%==============================================================================
\section{The Time-Count Identity}
\label{sec:time-count}
%==============================================================================

\subsection{Time as Generated, Not Given}

Newtonian and Lagrangian mechanics treat time as a parameter $t \in \Real$, an axis along which dynamical evolution is laid out. The state $x(t)$ is conceived as a curve indexed by this parameter. Time-reversibility, in this picture, is the statement that the curve can be traced in either direction.

We adopt a different view, motivated by the structure of bounded phase-space dynamics \citep{sachikonye2025temporal}. In any bounded oscillatory system, the elapsed time can be determined by counting cycles of a reference oscillator:
\begin{equation}
    t = \frac{\Mcount}{f},
    \label{eq:tmf}
\end{equation}
where $f$ is the oscillator frequency and $\Mcount$ is the number of cycles completed.

\begin{theorem}[Time-Count Identity]
\label{thm:time-count}
For a bounded oscillatory system with reference frequency $f$, the relation $t = \Mcount/f$ is not a measurement equation but a definition. Time is the partition count, in time-units. Equivalently, the operator equation
\begin{equation}
    \frac{\dM}{\dt} = f
\end{equation}
holds exactly, with no rounding or approximation: $\Mcount(t)$ is the integer number of cycles in $[0, t]$, and the rate at which $\Mcount$ increments \emph{is} the oscillator frequency.
\end{theorem}

\begin{proof}
The oscillator generates partition transitions at rate $f$ by definition of frequency: in time $t$, it completes $\Mcount = ft$ cycles. The relation is exact whenever ``time'' is operationally defined---whenever there is a procedure for measuring elapsed time. Every such procedure, in physics, reduces to counting cycles of some reference oscillator. The relation $t = \Mcount/f$ therefore states the operational definition of time itself.

The mathematical apparatus of Newtonian time as a continuous parameter is a limit of this counting procedure as $f \to \infty$. The continuous-time formalism is a useful idealisation, but it should not be confused with a description of how time exists physically.
\qed
\end{proof}

\subsection{Two Operations: Keeping and Telling}

The identity \eqref{eq:tmf} draws a sharp distinction between two physical operations that both involve time \citep{sachikonye2025temporal}:

\begin{itemize}[leftmargin=1.2em]
    \item \emph{Keeping time}: generating partition transitions at a known rate. Atomic clocks, quartz oscillators, and pendulum clocks keep time. The keeping is a physical process.

    \item \emph{Telling time}: reading the accumulated count from a partition record. Sundials, tree rings, and mass spectrometers tell time. The telling is a decoding operation on a frozen record.
\end{itemize}

Both yield the same number---elapsed time---but from different physical substrates. A clock that has stopped continues to tell the time at which it stopped, even though it has ceased to keep time. A mass spectrum, once recorded, tells the duration of the ion's journey indefinitely, long after the journey is over.

\subsection{The Implication for Time-Reversal}

The Time-Count Identity has an immediate consequence: any operation that purports to reverse time must reduce the partition count $\Mcount$. There is no separate ``time axis'' along which $t$ could be made to flow backward independently of $\Mcount$. To make time go backward is to make $\Mcount$ decrement.

We shall return to this implication repeatedly. The remainder of the paper is largely an unfolding of its consequences.

\subsection{The Fundamental Identity}

Combining the Time-Count Identity with the dynamical structure of bounded oscillatory motion yields the fundamental identity \citep{sachikonye2025counting}:
\begin{equation}
    \frac{\dM}{\dt} = \frac{\omega}{2\pi} = \frac{1}{\tp},
    \label{eq:fundamental}
\end{equation}
where $\omega$ is the oscillation frequency in radians per second and $\tp$ is the mean partition residence time. The three equalities express:
\begin{enumerate}[leftmargin=1.2em,itemsep=0.2em]
    \item The counting rate equals the oscillator frequency (definition of counting).
    \item Each angular cycle generates exactly one partition transition (geometric closure).
    \item The system vacates each partition state at the oscillator rate (residence-time identity).
\end{enumerate}

These are not three approximations; they are three coordinatisations of the same physical fact. Time, frequency, and partition counting are unit-conversions of a single quantity.

\begin{figure*}[!htbp]
    \centering
    \includegraphics[width=\textwidth]{figures/panel_E1_time_count_identity.png}
    \caption{\textbf{Time-Count Identity: $t = M/f$ verified across ten orders of magnitude in time and twelve orders of magnitude in frequency.}
    (\textbf{A}) Cycle count $M = ft$ as a function of elapsed time $t$ on logarithmic axes for six oscillator frequencies spanning $f \in [10^3, 10^{15}]$~Hz (kHz to PHz, viridis colour gradient). All six curves are strictly linear with unit slope, confirming that $M$ is exactly proportional to both $t$ and $f$. Markers are filled circles at the eleven test times $t \in [10^{-9}, 10^1]$~s; lines connect successive measurements within each frequency series. The strict linearity establishes the operational definition of time as a partition count.
    (\textbf{B}) Reconstruction error $|t_{\mathrm{rec}} - t|/t$ where $t_{\mathrm{rec}} = M/f$ is the time recovered from the count alone, plotted on semi-logarithmic axes. The dashed red line marks the floating-point precision floor at $10^{-15}$. All recoverable measurements (those with $ft \geq 1$) fall at or below this floor, demonstrating that the Time-Count Identity holds to numerical exactness rather than as an approximation. Sub-cycle conditions ($ft < 1$, omitted as NaN) correspond to the regime in which counting cannot resolve fractional cycles.
    (\textbf{C}) Three-dimensional surface of $\log_{10} M(\log_{10} t, \log_{10} f) = \log_{10} t + \log_{10} f$, plotted on a $30 \times 30$ grid covering $f \in [10^3, 10^{15}]$~Hz and $t \in [10^{-9}, 10^1]$~s. The surface is a perfect plane in log-coordinates, confirming the bilinear structure $M = ft$. The viridis colourmap encodes $\log_{10} M$, ranging from $\sim -6$ at the closest corner to $\sim +16$ at the farthest. The viewing angle (elevation $22^\circ$, azimuth $-50^\circ$) maximises visibility of the planar geometry.
    (\textbf{D}) Distribution of $\log_{10}$ relative reconstruction errors across all $66$ test conditions in which the count is resolvable ($ft \geq 1$). The histogram is concentrated at the floating-point floor near $\log_{10}(\mathrm{rel\,err}) \approx -16$, with no measurements above $-10$, confirming that the Time-Count Identity holds to machine precision uniformly across the tested parameter range.}
    \label{fig:panel-E1}
\end{figure*}

%==============================================================================
\section{The Loop Identity Lemma}
\label{sec:loop}
%==============================================================================

\subsection{Setup}

Consider a deterministic, formally invertible evolution operator $f: \StateSpace \to \StateSpace$ acting on a state space $\StateSpace$. Suppose three states $A, B, C \in \StateSpace$ satisfy:
\begin{equation}
    f(A) = B, \quad f(B) = C,
    \label{eq:abc}
\end{equation}
with $A \neq C$. The forward trajectory is therefore $A \to B \to C$. Reversibility implies:
\begin{equation}
    f^{-1}(B) = A, \quad f^{-1}(C) = B.
    \label{eq:abc-reverse}
\end{equation}

We now consider the loop $A \to B \to C \to B \to A \to B$, executed by alternating $f$ and $f^{-1}$. The question is: at the end of this loop, are we at the original $B$, or at a different ``$B$''?

\subsection{The Dilemma}

\begin{lemma}[Loop Identity Lemma]
\label{lem:loop}
For the loop $A \to B \to C \to B \to A \to B$, one of the following must hold:
\begin{enumerate}[leftmargin=1.5em,itemsep=0.2em]
    \item[(i)] $A = C$, contradicting the hypothesis $A \neq C$.
    \item[(ii)] The state $B$ encountered at different points in the loop is not a single object; it must carry partition structure beyond its $\StateSpace$-coordinates that distinguishes ``$B$ from $A$'' from ``$B$ from $C$.''
\end{enumerate}
\end{lemma}

\begin{proof}
Suppose, for contradiction, that neither (i) nor (ii) holds. Then $A \neq C$ and the state $B$ is fully characterised by its $\StateSpace$-coordinates alone (no additional partition structure).

The forward step $f$ at $B$ uniquely produces $C$ (determinism). The backward step $f^{-1}$ at $B$ uniquely produces $A$ (reversibility). If $B$ is unique (no partition structure), then both maps are well-defined functions of $B$ alone, yielding $C$ and $A$ respectively. So far, no contradiction.

Now consider the round trip $B \xrightarrow{f^{-1}} A \xrightarrow{f} B'$. By the assumption that $B$ is unique and the dynamics is deterministic, $B' = B$. From $B'$, applying $f$ yields $C$ (since $B' = B$ and $f(B) = C$).

Now consider the alternative round trip $B \xrightarrow{f} C \xrightarrow{f^{-1}} B''$. Again, $B'' = B$. From $B''$, applying $f^{-1}$ yields $A$.

So we have a single state $B$ with two distinguished neighbours: $C$ (in the forward direction) and $A$ (in the backward direction). The state $B$ must therefore encode both ``my forward continuation is $C$'' and ``my backward predecessor is $A$''. These are different pieces of information about the same state.

Now the crucial step: consider an observer who arrives at $B$ via different routes. Observer 1 arrives via $A \to B$. Observer 2 arrives via $C \to B$ (i.e., reverse-evolves $C$). If $B$ has no partition structure, Observer 1 and Observer 2 see exactly the same state. Both must therefore predict the same forward continuation. Observer 1, having arrived from $A$, predicts $f(B) = C$. Observer 2, having arrived from $C$, also predicts $f(B) = C$ (since $B$ is the same). But Observer 2's recent trajectory was $C \to B$; if forward and backward are not distinguished by the state itself, then ``forward'' for Observer 2 should be the continuation of $C \to B$, which is $A$ (the next reverse step). The two observers therefore predict different forward continuations from the same state. Contradiction.

The contradiction is resolved only if either $A = C$ (so both observers predict the same thing trivially), or $B$ carries partition structure that distinguishes the observers' contexts. \qed
\end{proof}

\subsection{The Operational Reading}

Lemma~\ref{lem:loop} states a structural fact about reversible deterministic systems: their states cannot be ``bare phase points'' if the dynamics is non-trivial. The state must carry historical or contextual information beyond its phase-space coordinates---which is exactly partition structure.

This observation is purely formal; it does not yet invoke physics. Its consequence is profound: any physical system whose observable state is identified with phase-space coordinates alone cannot also be reversibly evolving and non-trivially dynamic. At least one of these three properties must fail. In the partition framework, the failure is identified: \emph{observable state $\neq$ microstate}, in agreement with the framework of \citep{sachikonye2025loschmidt}.

\subsection{Application to Mass Spectrometry}

In mass spectrometry, the ion's microstate is given by its phase-space coordinates $(\mathbf{r}, \mathbf{p})$. The observable state at the detector is given by $(m/z, t_{\mathrm{arrival}}, I)$. The loop lemma forbids these two from coinciding: the observable state must encode partition structure---which physical analyzers were traversed, in what order, with what selection partitions---beyond what can be read from the ion's phase-space coordinates.

This formal observation has a direct physical content: \emph{the spectrum is not determined by the ion alone; it is determined by the ion together with the analyzer chain that processed it.} The reproducibility of mass spectra depends on both ion and instrument; the observable is irreducibly relational.

%==============================================================================
\section{The Sliding-Endpoint Theorem}
\label{sec:sliding}
%==============================================================================

\subsection{The Counterfactual}

Consider a mass-spectrometric measurement in which an ion is detected at time $t_C$. The observable record at the detector is
\begin{equation}
    O = \left(\Mcount(t_C), \, m/z, \, t_{\mathrm{arrival}}, \, I\right),
\end{equation}
where $\Mcount(t_C)$ is the partition count accumulated up to detection. By the Time-Count Identity, $\Mcount(t_C) = f t_C$. By the State-Mass Correspondence \citep{sachikonye2025counting}, $\Mcount(t_C) \leftrightarrow m/z$.

Suppose, hypothetically, that the dynamics were reversible in the strong sense: time-stretches could be deleted by running the inverse evolution. Then we should be able to construct a counterfactual: \emph{what would the same ion have produced as observable, had we detected it $\Delta t$ earlier?}

\subsection{Two Incompatible Answers}

Under the strong-reversibility hypothesis, two routes give different answers to the counterfactual.

\textbf{Route 1 (reversibility-preserves-identity).} The ion's physical identity (its mass, charge, internal state) is determined by its phase-space content. Detecting the same ion earlier means intercepting it before it reached the detector. Its identity at the earlier moment is the same as at the later moment, so the readout it would produce must encode the same mass: same $m/z$.

\textbf{Route 2 (counting-determines-readout).} By the Time-Count Identity, the count at $t_C - \Delta t$ is strictly less than the count at $t_C$ by $\Delta \Mcount = f \Delta t$. By State-Mass Correspondence, this corresponds to a different $m/z$ value: a different mass.

\textbf{Same ion, two different masses, depending on a freely-chosen stopping time.} This is the absurdity. Mass becomes a function of measurement bookkeeping rather than a property of the particle.

\subsection{The Theorem}

\begin{theorem}[Sliding-Endpoint Theorem]
\label{thm:sliding}
The empirical reproducibility of mass-spectrometric measurements is logically equivalent to the irreversibility of partition transitions. Equivalently:
\begin{equation}
    \text{MS reproducible} \;\Longleftrightarrow\; \text{partitions cannot be deleted}.
\end{equation}
\end{theorem}

\begin{proof}
($\Rightarrow$) Suppose MS is reproducible: the same ion under the same conditions yields the same $m/z$. Suppose, for contradiction, that partition transitions could be deleted---that is, that for some time-stretch $[t_C - \Delta t, t_C]$ along the ion's trajectory, the partitions actualised in that interval could be removed. By the Time-Count Identity, this corresponds to reducing the count from $\Mcount(t_C)$ to $\Mcount(t_C) - f\Delta t$. By State-Mass Correspondence, the readout would change accordingly. But by hypothesis the same ion yields the same readout, and the only way for the readout to be invariant under count-reduction is for the count itself to be invariant---which contradicts the assumption that the partitions were deleted. Therefore partitions cannot be deleted.

($\Leftarrow$) Suppose partitions cannot be deleted. Then the count $\Mcount(t_C)$ accumulated during an ion's trajectory is a well-defined function of the trajectory, and the readout determined by this count is a well-defined function of the ion and the chain. Two measurements of the same ion under the same conditions traverse the same partition structure, accumulate the same count, and yield the same readout. This is reproducibility. \qed
\end{proof}

\subsection{The Empirical Significance}

Theorem~\ref{thm:sliding} elevates the most ordinary fact about mass spectrometry---that it works---into an empirical proof of partition irreversibility. Every reproducible mass spectrum is a verification of the Second Law in operational form. There is no instrumental fact in physics that more directly confirms the irreversibility of nature.

The conventional rhetorical move in expositions of the Second Law is to argue that irreversibility is somehow surprising or in need of explanation. The Sliding-Endpoint Theorem reverses this: irreversibility is precisely the condition under which our most precise measurement instruments give reproducible results. The surprising thing would be if it were otherwise.

\begin{figure*}[!htbp]
    \centering
    \includegraphics[width=\textwidth]{figures/panel_E2_sliding_endpoint.png}
    \caption{\textbf{Sliding-Endpoint Theorem: MS reproducibility is logically equivalent to partition irreversibility.}
    (\textbf{A}) Reproducibility test for twenty representative ions distributed across $m/z \in [100, 1500]$. For each ion, five replicate measurements were performed under identical conditions (same scan duration, same trap field, same starting kinetic energy). Each ion's five $m/z$ readouts are overlaid as blue points; the scatter is invisible at the resolution of the plot, with replicate spread bounded by floating-point precision ($\sigma < 1.14 \times 10^{-16}$). This confirms that under conserved operating conditions, the partition count---and hence the $m/z$ readout---is a deterministic function of the ion and the chain.
    (\textbf{B}) Cycle count $M$ accumulated up to detection as a function of stop time $t_C$ across eleven values $t_C \in [0.5, 1.5]$~ms, for a representative ion at $m/z = 500$ in a $10$~MHz reference oscillator. The relationship is strictly linear and monotone-increasing with slope $f = 10^7$~Hz; no two stop times produce the same count. The strict monotonicity is the operational content of partition accumulation: any change in the stop time corresponds to a distinct partition record.
    (\textbf{C}) Three-dimensional surface of the hypothetical mass shift $\delta(m/z)$ that would result from deleting $\Delta M$ partitions from the count, plotted across $m/z \in [100, 1500]$ and $\Delta M \in [10^0, 10^4]$ on a $25 \times 25$ grid. The magma colourmap encodes the predicted shift, ranging from sub-Dalton to $\sim 50$~Da at the largest deletion. The surface visualises the empirically forbidden mass shifts: were partition deletion possible, every point on the surface would be a measurable consequence; the absence of such shifts in real measurements is the empirical signature of partition irreversibility.
    (\textbf{D}) Distribution of replicate spreads (standard deviation of $m/z$ across five replicates) for one hundred ions across the full $m/z$ range, on logarithmic horizontal axis. All spreads are bounded by floating-point precision (dashed red line at $\log_{10} = -15$), with the histogram concentrated at the precision floor. The complete absence of measurable spread is the operational manifestation of the Sliding-Endpoint Theorem: identical conditions yield identical readouts, demonstrating that the count is well-defined.}
    \label{fig:panel-E2}
\end{figure*}


%==============================================================================
\section{Constitutive Asymmetry of Operations}
\label{sec:constitutive}
%==============================================================================

\subsection{Trajectories Are Not Curves}

The arguments of Sections~\ref{sec:loop} and~\ref{sec:sliding} live within the framework of dynamical systems: state spaces, evolution operators, trajectories as parameterised curves. Within this framework, partition irreversibility is a property of trajectories: they cannot be retraced backward.

We now move to a deeper level of analysis. Real physical processes are not curves in continuous phase space. They are sequences of physical operations, each constitutively one-way.

\begin{definition}[Operation]
A physical operation $\Operation: \StateSpace \to \StateSpace$ is a partial map between observable states implemented by a specific physical mechanism. The mechanism is the means by which the operation is realised in the laboratory or in nature; different mechanisms producing the same input-output map are different operations.
\end{definition}

The distinction between an operation and its mathematical action is essential. A unitary gate in a quantum computer and a thermal annealing process can produce the same input-output map, but they are different operations realised by different mechanisms.

\subsection{One-Way Operations}

\begin{definition}[Constitutively One-Way Operation]
An operation $\Operation: \StateSpace \to \StateSpace$ is \emph{constitutively one-way} if no operation $\Operation^*$ exists in physical reality such that $\Operation^* \circ \Operation = \mathrm{id}_{\StateSpace}$ as physical processes (not merely as mathematical maps).
\end{definition}

The empirical world is full of constitutively one-way operations. A few examples:

\begin{itemize}[leftmargin=1.2em,itemsep=0.2em]
    \item \emph{Fertilisation:} the fusion of male and female gametes to form a zygote. There is no biological mechanism called ``un-fertilisation.'' Reverse processes (e.g., parthenogenesis induction) are different operations with different mechanisms.

    \item \emph{Combustion:} the oxidation of a fuel. There is no mechanism that reverses combustion as a single operation. Reduction via reaction with hydrogen, photosynthesis, and electrolysis exist, but each is a different operation with different intermediates.

    \item \emph{Mitosis:} the division of one cell into two. ``Cell fusion'' exists as a process, but it is not mitosis-run-backward; it has different mechanisms (membrane fusion proteins, fusogen activity) and different intermediates.

    \item \emph{Crystallisation:} the formation of an ordered solid. Dissolution is a different operation with different kinetic pathways.

    \item \emph{Detection (in MS):} the conversion of an ion impact into an electrical signal via secondary-electron cascade. ``Un-detection'' is not a process.
\end{itemize}

\subsection{The Maize Seed: A Case Study}

Consider a maize seed in your hand. Asking ``how was this seed produced?'' is a coherent biological question with a determinate answer: the seed came to be through a specific sequence of operations---pollination, fertilisation of the egg cell and central cell, formation of zygote and endosperm, mitotic divisions, cellular differentiation, accumulation of storage starch, dehydration, dormancy.

Asking ``can this seed be reached by reversing time?'' is, in contrast, an incoherent question. Reversing time would have to reverse each of the operations that produced it. But fertilisation has no reverse mechanism; mitosis has no reverse mechanism; differentiation has no reverse mechanism; dehydration has a reverse (rehydration) but it is a different operation, not dehydration-run-backward. The ``reversed trajectory'' that would yield the seed by going backward simply does not exist as a physical possibility---not because it is improbable, but because the operations it would consist of are not operations in the physical world.

The seed is the end-product of a sequence of constitutively one-way operations. Its existence in your hand is testimony to the directionality of those operations.

\subsection{The Mass Spectrometer Chain}

The same structural observation applies to the mass spectrometer.

\begin{table}[H]
\centering
\caption{Constitutively one-way operations in MS.}
\label{tab:ms-operations}
\begin{tabular}{p{0.4\columnwidth}p{0.5\columnwidth}}
\toprule
\textbf{Forward operation} & \textbf{``Reverse''} \\
\midrule
Ionisation (electron impact, electrospray, MALDI) & No mechanism. Recombination is a different process with different intermediates. \\
\addlinespace
Transmission through ion optics & Particle does not ``un-transmit''; focusing fields are intrinsically vectorial. \\
\addlinespace
Mass selection (quadrupole, magnetic sector) & The selection partition $\{\text{passes}, \neg\text{passes}\}$ commits the ion's identity to one branch. \\
\addlinespace
Fragmentation (CID, HCD, ETD) & Reassembly via collisional cooling exists, but is a different operation with different intermediates. \\
\addlinespace
Detection (secondary-electron cascade) & The multiplier amplifies; ``un-amplification'' is not a process. \\
\bottomrule
\end{tabular}
\end{table}

The MS journey is a sequence of operations, each constitutively one-way. The detection event is the end-product of this sequence, just as the maize seed is the end-product of its sequence of biological operations.

\subsection{The Theorem}

\begin{theorem}[Constitutive Asymmetry]
\label{thm:constitutive}
Every physical trajectory is a sequence of operations $\{\Operation_1, \Operation_2, \ldots, \Operation_k\}$, each of which is constitutively one-way. The composite trajectory has no reverse trajectory in the operational sense: no sequence of operations exists that would undo the forward trajectory as a physical process.
\end{theorem}

\begin{proof}[Proof sketch]
The argument proceeds by exhausting the candidates for ``reverse trajectory.'' Suppose, for contradiction, that for some trajectory $\Operation_k \circ \cdots \circ \Operation_1$, a reverse trajectory $\widetilde{\Operation}_1 \circ \cdots \circ \widetilde{\Operation}_k$ exists such that the composite is the identity. Then for each $i$, $\widetilde{\Operation}_i$ inverts $\Operation_i$ as a physical process. We consider three cases.

\emph{Case 1: $\widetilde{\Operation}_i$ does not exist.} Then no reverse operation is available. Contradiction with the assumption.

\emph{Case 2: $\widetilde{\Operation}_i$ exists but is a different operation.} Then $\widetilde{\Operation}_i$ has different mechanism, different intermediates, different thermodynamic signature than $\Operation_i$. As a process, it is not the inverse of $\Operation_i$ but a distinct process whose mathematical action happens to match $\Operation_i^{-1}$. The composite $\widetilde{\Operation}_i \circ \Operation_i$ is therefore a sequence of two distinct operations, not the identity-as-a-process.

\emph{Case 3: $\widetilde{\Operation}_i$ exists and is the same operation.} Then $\Operation_i$ is its own inverse. This is possible only if $\Operation_i$ is a symmetry transformation (e.g., a parity flip, a 180-degree rotation) that does not advance the partition structure. Such operations exist but do not constitute the bulk of physical trajectories. Generic operations (mass selection, ionisation, detection, mitosis, fertilisation) are not their own inverses.

In all generic cases, the reverse trajectory does not exist as a physical process. \qed
\end{proof}

\begin{corollary}[Reversibility Is a Coarse-Graining Artifact]
\label{cor:coarse-graining}
The time-reversibility of the continuous-flow description $\dot{x} = f(x,t)$ is an artifact of replacing operation-sequences with continuous flows. The flow description averages over the operational structure, and the averaging happens to be invariant under $t \to -t$. The averaging itself is not a physical process.
\end{corollary}

\subsection{The Implication for Loschmidt}

Theorem~\ref{thm:constitutive} undermines the very framing of Loschmidt's paradox. Loschmidt argued that since the equations of motion are time-reversible, reversed evolutions exist as solutions and we should occasionally observe them. The Constitutive Asymmetry Theorem responds: the equations of motion describe averages over operations, not the operations themselves. The reversibility of the equations does not correspond to the existence of reverse operations. Reverse operations are absent at the level of mechanisms, not merely improbable at the level of solutions.

%==============================================================================
\section{Rewind-as-Forward and the Incoherence of Time-Reversal}
\label{sec:rewind}
%==============================================================================

\subsection{The Movie Analogy}

Consider a film projector showing a movie. Forward play: frames 1, 2, 3, $\ldots$, $n$ are displayed in order. Rewind: frames $n, n-1, \ldots, 1$ are displayed in reverse order.

Now ask: during the rewind, does time flow backward? Of course not. The rewinding takes time. The projector mechanism advances forward in time as it rewinds the film. Photons are emitted forward; the audience perceives them forward; the operator's wristwatch ticks forward. The frames go in reverse order, but the reverse-ordering is itself a forward physical process.

This observation is more than an analogy.

\begin{principle}[Rewind-as-Forward]
\label{princ:rewind}
A ``rewound'' sequence is not the original sequence run backward. It is a new forward sequence whose state-content is the reverse-ordered list of states from the original. The rewinding itself is a forward operation that increments the global partition count.
\end{principle}

The principle applies to every claim of ``time-reversal'' in physics. Maxwell's demon reversing molecular velocities? The demon must act forward in time, perform forward measurements, expend forward thermodynamic resources. The ``reversed'' molecular trajectory that ensues is a new forward trajectory whose state-content happens to mimic the reverse-ordering of the original states. By the Constitutive Asymmetry Theorem, the operations constituting this new trajectory are not the inverses of the original operations.

\subsection{Solution Sets vs. Physical Processes}

The mathematical operation $T: x(t) \to x(-t)$ is well-defined as a map on the solution set of the equations of motion. It takes a solution and produces another solution. As an automorphism of the solution set, it is no more problematic than the map ``send $f(x)$ to $f(-x)$'' on a function space.

But this mathematical automorphism does not correspond to a physical operation on a system. There is no physical procedure whose enactment maps the system from solution $x(t)$ to solution $x(-t)$. The mathematical operation acts on representations; physical operations act on systems.

The conflation of these two levels---automorphisms of solution sets and physical operations on systems---is the source of Loschmidt's paradox. The paradox arises by treating the mathematical automorphism as if it were a physical operation, then asking why we don't observe its consequences.

\subsection{The Incoherence Theorem}

\begin{theorem}[Incoherence of Time-Reversal]
\label{thm:incoherence}
There is no physical operation whose enactment decrements the global partition count. The mathematical operation $T: x(t) \to x(-t)$ is an automorphism of the solution set of the equations of motion, not a physical process. Time-reversibility is therefore not a property of physical dynamics but a relation on solution representations. The question ``is time-reversal possible in physical systems?'' is a category error.
\end{theorem}

\begin{proof}
Suppose, for contradiction, that some physical operation $\Operation^*$ decrements the partition count: $\Mcount$ after $\Operation^*$ is less than $\Mcount$ before $\Operation^*$. The enactment of $\Operation^*$ is itself a physical process, occurring during some time-interval $\Delta t > 0$. By the Time-Count Identity, the partition count $\Mcount$ generated during $\Delta t$ is at least $f \Delta t > 0$, where $f$ is the system's reference oscillator frequency. Therefore the count after $\Operation^*$ exceeds the count before $\Operation^*$ by at least $f \Delta t$. This contradicts the assumption that $\Operation^*$ decrements the count.

The mathematical operation $T: x(t) \to x(-t)$ acts on solution functions, not on physical systems. Its application is an act of mathematical reasoning, not a laboratory procedure; no time elapses in the system being represented when we apply $T$. The mathematical operation has no operational analogue: there is no laboratory procedure whose execution corresponds to applying $T$ to a system.

Time-reversibility, as a property of the solution set under $T$, is therefore a feature of representations, not of physical processes. Asking whether physical dynamics is ``time-reversible'' is asking whether a relation on representations holds of physical reality---a question that has the form of a category error. \qed
\end{proof}

\begin{theorem}[Refined Statement: Inverses Cannot Decrement the Count]
\label{thm:refined-inverse}
For any operation $\Operation$ with mathematical inverse $\Operation^{-1}$ existing as a physical operation, the application of $\Operation^{-1}$ is itself an operation that increments the partition count by at least one. Therefore $\Operation^{-1}$ does not undo $\Operation$ at the partition-count level: $\Mcount(\Operation^{-1} \circ \Operation) > \Mcount(\mathrm{identity})$.
\end{theorem}

\begin{proof}
The application of $\Operation^{-1}$ is itself a physical process consuming some time $\Delta t > 0$. During this time, the system's reference oscillator generates at least $f \Delta t \geq 1$ partition transitions. The state may return to its original phase-space coordinates after $\Operation^{-1} \circ \Operation$, but the partition count is strictly larger than before $\Operation$ was applied. \qed
\end{proof}

Theorem~\ref{thm:refined-inverse} catches the right level of nuance. It does not deny that some operations have mathematical inverses; it observes that even when they do, the application of the inverse is itself a forward partition event. The phase-space trajectory may close on itself; the partition count cannot.

\subsection{The Resolution of Loschmidt}

We now state the resolution that the Incoherence Theorem provides. Loschmidt's paradox asked: if the equations of motion are time-reversible, why is the macroscopic world irreversible?

The resolution: the question is malformed. The equations of motion describe automorphisms of solution sets. Macroscopic irreversibility describes properties of physical processes. These are different categories. The question ``why is one irreversible if the other is reversible?'' presupposes that the two are commensurable. They are not.

Macroscopic processes are sequences of operations, each constitutively one-way (Theorem~\ref{thm:constitutive}). Their irreversibility is built into what they are. The mathematical automorphisms of the solution set are features of representations and have no operational realisation. The two never compete; they exist at different levels.

The Second Law of Thermodynamics, in this view, is not a statistical regularity but the operational signature of the constitutive directionality of physical operations.

\subsection{The Asymmetry of Becoming}
\label{sec:becoming}

A subtle but decisive aspect of the incoherence of time-reversal concerns the structure of \emph{becoming}---the operational mode by which a state is reached.

\begin{principle}[Becoming Is Forward-Only]
\label{princ:becoming}
A state can be the \emph{target of becoming} only if it lies in the future of the current state. Past states are not future-targets to be reached; they are way-stations already-passed-through. The two modes of existence---``target-of-becoming'' and ``way-station-passed-through''---are categorically distinct, not merely temporally distinct.
\end{principle}

The principle has direct consequences for the Loschmidt setup. To ``reverse'' an ion to a past state is to ask the ion to \emph{become} what it was at that past time. But the past version of the ion existed precisely because it was not yet the future version; its past identity is constituted, in part, by the not-yet-having-become of every subsequent state. Asking the ion now (in its future-state) to become its past-state is asking it to undo the becoming that constituted the past-state's identity. There is no operation that performs such an unbecoming, because the past-state's identity was never available as a future-target; it was always a way-station.

\begin{theorem}[No Path Backward]
\label{thm:no-path-backward}
Let $S_0$ denote the state of a system at time $t_0$, and let $S_1$ denote the state at time $t_1 > t_0$. There exists no sequence of operations $\{\Operation_1, \ldots, \Operation_k\}$ whose enactment from $S_1$ produces $S_0$ as a target-of-becoming.
\end{theorem}

\begin{proof}
Suppose, for contradiction, that such a sequence exists. Then $S_0$ is a target-of-becoming reached from $S_1$ via the operations $\{\Operation_i\}$. By Principle~\ref{princ:becoming}, a target-of-becoming lies in the future of the current state. Therefore $S_0$ lies in the future of $S_1$. But by hypothesis $S_0$ occurs at $t_0 < t_1$, so $S_0$ lies in the past of $S_1$. A state cannot lie both in the past and the future of $S_1$. Contradiction.

The contradiction is not avoidable by relabeling: the past-state's identity is constituted by its not-yet-having-become the subsequent states. To re-instantiate $S_0$ at a later time $t_2 > t_1$ would produce a state $S_0^{(2)}$ that is not $S_0$ but a configurationally similar later-state---one whose identity is constituted by an entirely different becoming-history. \qed
\end{proof}

\begin{corollary}[Past States Are Not Future-Targets]
\label{cor:past-not-future}
The Loschmidt operation ``reverse all velocities and let the system evolve'' does not produce the past state. Even when the resulting trajectory traverses configurations spatially identical to the original past, the produced states are not the past states. They are new future-states whose identity is constituted by the velocity-reversal-and-subsequent-evolution becoming-history, not by the original initial-condition-and-evolution becoming-history.
\end{corollary}

The Asymmetry of Becoming is the deepest reason why time-reversal is incoherent. Earlier sections established that no physical operation decrements the partition count. The present principle adds: even if such an operation were imagined, the state it produced would not be the past state---because the past state's identity is constituted by its not-being-future, and the produced state's identity is constituted by its having-been-arrived-at-from-the-future.

This places the final structural point in the framework. A point in state-space can only be reached \emph{forward}, because reaching is becoming, and becoming is forward-only. The past is not unreachable because of distance or difficulty; it is unreachable because it is not the kind of thing that can be reached.

\subsection{Oscillation as Forward Re-Placement}
\label{sec:oscillation-forward}

A natural objection arises at this point. Oscillatory and periodic systems---pendulums, atomic orbits, ion trajectories in mass spectrometers---appear to ``return'' to prior states. A pendulum at the bottom of its swing repeatedly traverses configurations that are nearly identical from cycle to cycle. Does this not constitute a counterexample to the forward-only character of operations? Does the pendulum not ``go back'' to where it was?

The objection is dissolved by recognising that oscillation is not return. It is forward re-placement.

\begin{principle}[Oscillation as Forward Re-Placement]
\label{princ:oscillation}
For any oscillatory or repetitive process, the apparent ``return'' to a prior configuration is in fact a forward placement of a configurationally-similar state in the future. The configurational similarity between cycle $n$ and cycle $n+1$ is genuine; the operational identity is not. Each ``return'' is a new becoming-event whose target lies in the future of the prior cycle, not in its past.
\end{principle}

Consider a one-dimensional harmonic oscillator $x(t) = A \cos(\omega t)$. At time $t_1 = 0$, the oscillator is at position $x_1 = A$. At time $t_2 = 2\pi/\omega$, the oscillator is again at $x_2 = A$. Naively, the oscillator has ``returned'' to its initial position. But the partition count has incremented by exactly one cycle: $\Mcount(t_2) - \Mcount(t_1) = 1$. The configurations $x_1$ and $x_2$ are spatially identical but operationally distinct.

\begin{theorem}[Cycle Distinguishability]
\label{thm:cycle-distinguishability}
Two states of an oscillatory system at configurationally-identical positions but at different cycle indices $n$ and $n+k$ ($k \geq 1$) are operationally distinct: their partition counts differ by exactly $k$, and their becoming-histories have lengths differing by $k$ cycles' worth of partition transitions.
\end{theorem}

\begin{proof}
By the Time-Count Identity, $\Mcount(t) = ft$ where $f = \omega/(2\pi)$. If $t_1$ and $t_2$ are configurationally-identical times satisfying $t_2 - t_1 = k \cdot 2\pi/\omega$, then $\Mcount(t_2) - \Mcount(t_1) = f \cdot 2\pi k / \omega = k$. The two states have the same phase-space coordinates but different partition counts. By Lemma~\ref{lem:loop} (Loop Identity), states with distinct partition counts are operationally distinct. \qed
\end{proof}

\begin{corollary}[The Oscillator's Cycle is a Forward Becoming]
\label{cor:oscillator-forward}
Each cycle of an oscillator is a fresh becoming-event placing the configurationally-prior state in the future. The cycle does not retrieve the past; it produces a new instance whose future-character is what makes it accessible at all.
\end{corollary}

The corollary explains, in operational terms, why oscillators function as time-keepers. If the oscillator's apparent ``return'' were a true return to the past state, the partition count would oscillate between two values rather than increment monotonically; mass spectrometry would not work; clocks would not keep time. The fact that clocks do keep time and mass spectrometers do produce reproducible spectra is the operational signature that oscillation is forward re-placement, not return.

This principle is the key to understanding why the partition framework reconciles periodic motion with constitutive asymmetry. The equations of motion for an oscillator admit periodic solutions; these solutions are mathematical objects on the solution set, invariant under temporal translation. The physical oscillator, however, does not realise periodicity by returning to past states. It realises periodicity by repeatedly placing similar configurations in the future. The partition count tracks this forward re-placement; the count, not the configuration, is the time.

A mass spectrometer's ion trap is a paradigm case. The ion oscillates within the trap, traversing configurationally similar positions thousands of times during a single scan. Each oscillation increments the partition count. The accumulated count at the moment of detection is the elapsed time of the ion's journey---and, by the Mass-Time-Identity-Count Equivalence (Theorem~\ref{thm:equivalence}), the mass of the ion. The mass spectrometer reads mass as accumulated forward re-placement of the ion's identity.

\begin{figure*}[!htbp]
    \centering
    \includegraphics[width=\textwidth]{figures/panel_E3_rewind_as_forward.png}
    \caption{\textbf{Rewind-as-Forward Principle: every reversal operation increments the partition count.}
    (\textbf{A}) Cumulative partition count along a five-stage forward analyzer chain (blue circles, stages 1--5) followed by a five-stage reverse chain (red squares, stages 6--10). The forward chain accumulates $M_{\mathrm{forward}} = 4{,}100$ cycles over stage durations $\{10, 50, 100, 200, 50\}$~$\mu$s; the reverse chain---which would, on the naive view, ``undo'' the forward---instead adds another $4{,}100$ cycles, yielding $M_{\mathrm{after\,reversal}} = 8{,}200$. The grey dashed line marks the forward maximum; the count never returns to or below this line. The plot is the operational content of the Rewind-as-Forward Principle: reversal is itself a forward operation.
    (\textbf{B}) Cycle increment $\Delta M$ produced by each of twenty randomly-generated inverse operations (durations $\in [10^{-7}, 10^{-5}]$~s drawn uniformly), shown as green bars on a logarithmic vertical axis. The dashed red horizontal line marks the theoretical minimum increment of one cycle ($k_B T \ln 2$ at the operational level). All twenty bars exceed this minimum, with the smallest increment being $1.0$ cycles ($100$~ns at the $10$~MHz reference). No inverse operation can produce a smaller increment because its enactment is itself a partition transition.
    (\textbf{C}) Three-dimensional scatter of thirty round-trip trials in $(M_{\mathrm{forward}}, M_{\mathrm{after\,reversal}}, \mathrm{ratio})$ space, where the ratio is $M_{\mathrm{after\,reversal}}/M_{\mathrm{forward}}$. Plasma colourmap encodes the ratio. All thirty points cluster near ratio $= 2.0$, with values bounded strictly above unity. Each round trip---a forward chain followed by a same-time reverse chain---produces an after-state with exactly twice the partition count of the before-state, a consequence of Theorem~\ref{thm:refined-inverse}. The viewing angle (elevation $20^\circ$, azimuth $-60^\circ$) emphasises the planar arrangement of points along the diagonal $M_{\mathrm{after}} = 2 M_{\mathrm{forward}}$.
    (\textbf{D}) Distribution of round-trip ratios $M_{\mathrm{after}}/M_{\mathrm{forward}}$ across thirty trials, shown as a purple histogram. The dashed red line marks the forbidden ratio of unity (where reversal would undo the forward operation); no trials fall in this region. The dashed green line marks the theoretical ratio of $2.0$ predicted by same-time round trips; the histogram is concentrated at this value. The empirical absence of ratios $\leq 1$ is the structural signature of partition irreversibility.}
    \label{fig:panel-E3}
\end{figure*}

%==============================================================================
\section{Identity as Fixed Point of Negation}
\label{sec:negation}
%==============================================================================

\subsection{The Persistence of a Particle}

Consider a particle $A$ at position $\mathbf{r}_1$ at time $t_1$. At a later time $t_2$, the particle is at position $\mathbf{r}_2$. We say it is ``the same particle.'' What does this mean?

In Newtonian physics, the answer is a continuity-of-trajectory statement: there exists a smooth path from $(\mathbf{r}_1, t_1)$ to $(\mathbf{r}_2, t_2)$ along which the particle has been continuously present. The identity of $A$ is preserved by the trajectory.

In partition theory, this answer is incomplete. The particle is at $\mathbf{r}_2$ at $t_2$ only because it has undergone the partition transitions between $t_1$ and $t_2$. Without those transitions, there is no $A$ at $t_2$ at all---there is nothing. The particle's persistence is not preserved by some abstract identity-through-time relation; it is \emph{constituted} by the sequence of partition transitions that the particle has undergone.

\subsection{The Operational Content of Identity}

Let $\Negation_t$ denote the negation operator at time $t$: the operator whose action on a state-space element $x$ produces ``$x$ if $x$ survives all alternative actualisations at $t$, otherwise undefined.''

\begin{theorem}[Identity as Fixed Point]
\label{thm:identity-fixed-point}
The persistence of particle $A$ from $t_1$ to $t_2$ is the property: $A$ is a fixed point of $\Negation_t$ for every $t \in [t_1, t_2]$. That is,
\begin{equation}
    \Negation_t(A) = A \quad \text{for all } t \in [t_1, t_2].
\end{equation}
Equivalently, $A$ persists if and only if at every moment, all alternative actualisations fail to displace it.
\end{theorem}

\begin{proof}
Following the negation framework of \citep{sachikonye2025partition,sachikonye2025loschmidt}, an actualisation occurs through the failure of competing negations. The persistence of $A$ from $t_1$ to $t_2$ is therefore the joint property: at every moment in this interval, $A$ continues to be the fixed point of the negation operator. If at some $t \in [t_1, t_2]$ there were a competing actualisation $A'$ that displaced $A$, then $A$ would not persist; the persistence assumption is therefore equivalent to the fixed-point property. \qed
\end{proof}

\subsection{Double Negation and the Forward Direction}

A particle at $t_2$ that ``is still $A$'' is, in the language of partition theory, a particle that is ``not not-$A$''---the double negation. But the double negation is not the trivial logical identity $\neg\neg A = A$. It is the operational statement that $A$ has survived the chain of partition transitions between $t_1$ and $t_2$.

Crucially, the chain of partition transitions has a direction. The fixed-point property holds along the forward time-direction. There is no symmetric statement ``$A$ at $t_1$ persists from $t_2$''---that is not a coherent claim, because persistence is operationally constituted by partition survival, and partition survival is a forward operation.

\begin{corollary}[Identity Inherits Constitutive Asymmetry]
\label{cor:identity-asymmetry}
The persistence of particle identity through time is constitutively asymmetric. It is a fixed-point property of forward partition operations, with no symmetric backward analogue.
\end{corollary}

This corollary is the conceptual bridge between the Constitutive Asymmetry Theorem and the partition framework's third route to mass.

\subsection{Mass as Survival Record}

The mass of a particle, in the partition framework, is the accumulated record of its persistence---specifically, the residue from partition lag during the particle's continuous self-actualisation \citep{sachikonye2025partition}. The particle, as a localised oscillatory mode, must continuously partition itself from the background to maintain its identity. Each self-partition cycle generates residue. The residue accumulates as mass.

\begin{theorem}[Mass as Accumulated Persistence]
\label{thm:mass-persistence}
The rest mass $m_0$ of a particle is the integrated record of its survival as a fixed point of negation:
\begin{equation}
    m_0 c^2 = \hbar \omega_0,
\end{equation}
where $\omega_0$ is the frequency of the particle's self-actualisation cycles.
\end{theorem}

The proof follows from the analysis of \citep{sachikonye2025partition} (Theorem 11.1) and is reproduced here for completeness in Section~\ref{sec:routes}.

%==============================================================================
\section{Source-Analyzer Indeterminacy}
\label{sec:source}
%==============================================================================

\subsection{Setup}

Consider an ion that traverses a chain of analyzer stages $A_1, A_2, \ldots, A_k$ before arriving at the detector $D$. At $D$, we record the observable $O = (m/z, t_{\mathrm{arrival}}, I)$.

The analyzer chain is a composition of partition operations. Each $A_i$ implements a selection partition $\{\text{passes}_i, \neg\text{passes}_i\}$; ions falling into $\neg\text{passes}_i$ are eliminated. The detected ion is the one for whom every selection succeeded:
\begin{equation}
    \text{ion} \in \bigcap_{i=1}^{k} \text{passes}_i.
\end{equation}

\subsection{The Indeterminacy Theorem}

\begin{theorem}[Source-Analyzer Indeterminacy]
\label{thm:source-indeterminacy}
The detector observable $O$ does not uniquely determine the analyzer chain $\{A_i\}$ that produced it. From $O$ alone, one can recover the fixed-point value of the composed negation operator $\Negation = \Negation_{A_k} \circ \cdots \circ \Negation_{A_1}$, but not the operator itself.
\end{theorem}

\begin{proof}
By the Loop Identity Lemma (Lemma~\ref{lem:loop}), the observable state at $D$ does not coincide with the ion's microstate; it must additionally encode partition structure. The relevant partition structure includes the chain of selection partitions $\{\text{passes}_i\}$.

Now consider a different chain $A'_1, A'_2, \ldots, A'_{k'}$ that, applied to a different initial ion, produces an ion arriving at $D$ with the same observable $O$. Such alternative chains exist: a TOF analyzer and a quadrupole analyzer can be configured to produce the same $m/z$ readout despite their different selection partitions. The distinguishing partition structure is the chain itself, encoded in the ion-analyzer correlation, not in the ion alone.

At the moment of detection, the ion-analyzer correlation is severed: the ion's content is committed to the detector's electrical signal, and the analyzer is freed to receive the next ion. The partition structure encoding the chain is therefore not present in the ion at $D$ in any way recoverable from $O$.

Consequently, the observable $O$ encodes only the fixed-point value (the mass) of the composed negation operator, not the operator's identity. The chain $\{A_i\}$ is undecidable from $O$ alone. \qed
\end{proof}

\subsection{The Operational Reading}

Theorem~\ref{thm:source-indeterminacy} has a striking operational consequence. The mass spectrometer, considered as an instrument, is partially blind to its own internal mechanism at the level of its output. Different chains can produce identical readouts. The output records the result of the chain, not the chain itself.

This is the operational content of Route III to mass: the ion is the fixed point of the negation operator $\Negation$, and the mass is the fixed-point value. The operator's identity is lost in the fixed point, just as the path of integration is lost in the value of a definite integral.

\subsection{Categorical Entropy and Source Information}

The information lost in the source indeterminacy has a quantitative measure. By the analysis of \citep{sachikonye2025counting}, each analyzer stage generates categorical entropy
\begin{equation}
    \Delta S_i \geq \kB \ln 2.
\end{equation}
The total categorical entropy generated along the chain is
\begin{equation}
    \Delta S_{\mathrm{chain}} = \sum_{i=1}^{k} \kB \ln\!\left(2 + \frac{|\delta\phi_i|}{100}\right) \geq k \kB \ln 2.
    \label{eq:chain-entropy}
\end{equation}
This entropy is precisely the quantity of source information committed to ion-analyzer correlations and made inaccessible at the detector. It is the operational measure of the chain's identity, lost when the ion arrives at $D$.

\begin{figure*}[!htbp]
    \centering
    \includegraphics[width=\textwidth]{figures/panel_E4_source_indeterminacy.png}
    \caption{\textbf{Source-Analyzer Indeterminacy: thirty distinct analyzer chains converge to identical detector readouts.}
    (\textbf{A}) Three-dimensional scatter of thirty analyzer chains in (number-of-stages, total partition count $M$, $m/z$ readout) space, with each chain colour-coded by its stage count (viridis, ranging from $2$ to $7$ stages). All thirty chains converge to the same total $M = 10^4$ cycles and the same $m/z = 500$~Da despite their differing internal structures. The convergence in two of the three coordinates---visible as the line of points at the front of the plot---demonstrates that the detector observable does not encode the chain's compositional structure. Stage durations within each chain were drawn from a Dirichlet distribution to enforce the total-time constraint while permitting maximal compositional diversity.
    (\textbf{B}) Histogram of the number of stages per chain across the thirty constructed chains, with bin width $1$ and stage counts ranging from $2$ to $7$. The distribution covers the full range with non-trivial population in each bin, demonstrating that the chains genuinely differ in their compositional cardinality, not merely in their stage durations. The diversity in stage count is the operational content of source indeterminacy: chains with $2$ stages and chains with $7$ stages produce structurally identical detector observables.
    (\textbf{C}) Distribution of source-information per chain in bits, computed as the Shannon entropy of the normalised stage-time distribution within each chain. The mean source-entropy is $3.54$~bits per chain (red dashed line), corresponding to approximately $3.4 \times 10^{-23}$~J/K of categorical entropy committed to the ion-analyzer correlations. This quantity is the average amount of source information that is committed to the analyzer hardware at the moment of detection, and that is irrecoverable from the detector readout alone.
    (\textbf{D}) Pairwise distinctness matrix for the thirty chains. Each cell $(i, j)$ encodes whether chains $i$ and $j$ are operationally distinct (white) or identical (black). The matrix is fully white off the diagonal: of $\binom{30}{2} = 435$ chain-pairs, all $435$ are pairwise distinct (different stage counts or different stage-duration distributions). The diagonal (self-comparisons) is by definition zero. The complete absence of off-diagonal black cells is the structural validation of the Source-Analyzer Indeterminacy Theorem: distinct operator chains genuinely produce identical fixed-point values.}
    \label{fig:panel-E4}
\end{figure*}

%==============================================================================
\section{The Three Routes to Mass and Their Equivalence}
\label{sec:routes}
%==============================================================================

\subsection{The Three Routes Recalled}

The partition framework offers three distinct routes to the concept of mass, each developed in detail in \citep{sachikonye2025partition}. We summarise:

\begin{itemize}[leftmargin=1.2em,itemsep=0.3em]
    \item \emph{Route I (Residue):} Mass is the accumulated residue from partition lag. A localised oscillatory mode, in maintaining its identity, generates residue at ratio $(b^d - 1)/b^d$ per cycle. The accumulated residue is the rest energy.

    \item \emph{Route II (Confinement):} Mass is the partition confinement cost. A localised mode confined to spatial extent $L$ has minimum oscillation frequency $\omega_0 > 0$, and the inertial mass is $m_0 = \hbar\omega_0/c^2$.

    \item \emph{Route III (Negation Fixed Point):} Mass is the spectrum of fixed points of the universal negation operator. A mode at frequency $\omega_0$ persists if and only if it survives all competing negations: $\Negation(\omega_0) = \omega_0$.
\end{itemize}

\subsection{Equivalence under the Partition-Counting Framework}

\begin{theorem}[Three-Route Equivalence]
\label{thm:three-routes}
The three routes to mass yield the same physical quantity. They are unit-conversions of a single mass-observable $m_0$.
\end{theorem}

\begin{proof}[Proof sketch]
\emph{Route I $\Leftrightarrow$ Route II.} The accumulated residue per cycle is $(b^d - 1)\hbar\omega_0/b^d$, and the actualised structure is $\hbar\omega_0/b^d$. The total rest energy is therefore $\hbar\omega_0$. By Route II, $m_0 c^2 = \hbar\omega_0$. Therefore Route I and Route II give the same mass value; they differ only in interpretation (memory vs. confinement cost).

\emph{Route II $\Leftrightarrow$ Route III.} The fixed points of the negation operator are the values $\omega_0$ at which the mode survives its own negation filter. By Route II, these correspond to specific masses $m_0 = \hbar\omega_0/c^2$. Therefore the mass spectrum is the fixed-point spectrum.

\emph{Route I $\Leftrightarrow$ Route III.} The accumulated residue is the operational record of all the negations that the mode survived. The mass-as-residue equals the mass-as-fixed-point of the negation operator, viewed as the sum of evidences of survival.

The three routes are therefore three coordinatisations of the same underlying quantity. \qed
\end{proof}

\subsection{Source Indeterminacy as the Operational Recovery of Route III}

The Source-Analyzer Indeterminacy Theorem (Theorem~\ref{thm:source-indeterminacy}) gives operational content to Route III. The ion that hits the detector is the fixed point of the chain's composed negation operator; the mass we read is the fixed-point value. The operator itself---the chain---is not recoverable from the readout. This is exactly what Route III claims: the mass spectrum is the spectrum of fixed-point values of the negation operator, with the operator's specific form being structurally invisible at the level of the fixed point.

The experimental verification of Route III is therefore not a separate measurement but a structural feature of every mass-spectrometric observation. Every spectrum is a Route-III measurement of its constituents.

\subsection{The Mass-Time-Identity-Count Equivalence}

Combining the Time-Count Identity (Theorem~\ref{thm:time-count}), Identity as Fixed Point (Theorem~\ref{thm:identity-fixed-point}), Mass as Accumulated Persistence (Theorem~\ref{thm:mass-persistence}), and the Three-Route Equivalence (Theorem~\ref{thm:three-routes}) yields:

\begin{theorem}[Mass-Time-Identity-Count Equivalence]
\label{thm:equivalence}
For a localised oscillatory mode (a particle), the following four physical quantities are equivalent up to unit-conversion factors:
\begin{enumerate}[leftmargin=1.5em,itemsep=0.2em]
    \item Mass $m$ (in kg).
    \item Elapsed proper time $\tau$ since the mode's emergence (in s).
    \item Identity (the cumulative survival record as fixed-point of negation).
    \item Partition count $\Mcount$ (dimensionless).
\end{enumerate}
The conversions are:
\begin{align}
    m &= \frac{\hbar \omega_0}{c^2}, \\
    \omega_0 &= 2\pi f, \\
    \Mcount &= f \tau, \\
    \tau &= \Mcount / f.
\end{align}
Identity is the property that $\Mcount$ counts in fact the partition cycles of one continuously-actualising mode; it is encoded in the dynamics rather than in a separate scalar.
\end{theorem}

The equivalence is not metaphysical speculation. It is the structural content of the partition framework, derived from the Time-Count Identity and the negation route to mass. Every mass measurement is, simultaneously, a measurement of the elapsed proper time of the particle's continuous self-actualisation, expressed in mass-units rather than time-units.

\begin{figure*}[!htbp]
    \centering
    \includegraphics[width=\textwidth]{figures/panel_E5_equivalence.png}
    \caption{\textbf{Mass-Time-Identity-Count Equivalence: mass, frequency, period, and partition count are unit-conversions of a single quantity across ten ions spanning four orders of magnitude in mass.}
    (\textbf{A}) Log-log plot of rest frequency $\omega_0 = mc^2/\hbar$ as a function of rest mass $m$ (kg) for ten test ions from proton ($m \approx 1.67 \times 10^{-27}$~kg) to lysozyme-10+ ($m \approx 2.38 \times 10^{-23}$~kg), drawn as blue circles connected by a line. The strict linearity with unit slope confirms that $\omega_0$ is exactly proportional to $m$, with proportionality constant $c^2/\hbar \approx 8.5 \times 10^{50}$~rad/(s$\cdot$kg). The linearity over four mass-decades demonstrates the framework's parameter-free unit-conversion structure.
    (\textbf{B}) Relative error in the energy-frequency identity $|E_{\mathrm{rest}} - \hbar\omega_0|/E_{\mathrm{rest}}$ for each ion, plotted on a semi-logarithmic vertical axis. The dashed red line marks floating-point precision at $10^{-15}$. All ten ions show errors at or near $10^{-17}$, far below the precision floor, confirming that the identity holds to machine precision across the entire mass range. The horizontal axis labels each ion by its abbreviated name (proton, He, C, Gly, Ala, Reserpine, Bradykinin, Substance~P, Insulin-5+, Lysozyme-10+).
    (\textbf{C}) Three-dimensional scatter of the ten ions in $(\log_{10} m, \log_{10} \omega_0, \log_{10} \tau_p)$ space, with viridis colourmap encoding $\log_{10} m$. Points lie on a strict line in the three-dimensional space, confirming the rigid unit-conversion relations $\omega_0 = m c^2/\hbar$ and $\tau_p = 2\pi/\omega_0$. The strict colinearity is the visual signature of the Mass-Time-Identity-Count Equivalence: any one of the four quantities determines the other three exactly.
    (\textbf{D}) Partition count $M = f_0 \cdot T_{\mathrm{obs}}$ accumulated by each ion in a $T_{\mathrm{obs}} = 1$~ms observation window, plotted on a semi-logarithmic vertical axis. Counts span from $\sim 10^{20}$ for protons to $\sim 10^{24}$ for lysozyme-10+, a five-order-of-magnitude range corresponding to the ions' rest-frequency range. The orange-marker series is monotone in mass, demonstrating that for fixed observation time, the partition count is a direct proxy for mass.}
    \label{fig:panel-E5}
\end{figure*}


\subsection{The Mass Spectrum as Visible Time}

A mass spectrum is therefore a record of elapsed times: each peak's position on the $m/z$ axis is, up to instrumental conversion factors, the elapsed proper time of an ion's self-actualisation cycle structure. The instrument does not measure mass and incidentally encode time; it tells time and incidentally encodes mass. The two are the same quantity.

Stated more provocatively: every mass spectrum is the arrow of time made visible. The spread of peaks across the $m/z$ axis is the spread of distinct elapsed times of the constituent particles. The reproducibility of peak positions is the Constitutive Asymmetry Theorem in operational form.

%==============================================================================
\section{Resolution of Loschmidt's Paradox}
\label{sec:resolution}
%==============================================================================

\subsection{The Resolution Stated}

We now state the resolution of Loschmidt's paradox.

\begin{theorem}[Resolution of Loschmidt's Paradox]
\label{thm:resolution}
Loschmidt's paradox arises from a category error: the conflation of automorphisms of solution sets with physical operations on systems. The mathematical operation $T: x(t) \to x(-t)$ is an automorphism of the solution set of the equations of motion; it is not a physical operation. Physical operations are constitutively one-way (Theorem~\ref{thm:constitutive}); their composite trajectories have no reverse trajectories in the operational sense; and any operation purporting to enact a reversal is itself a forward operation that increments the partition count (Theorem~\ref{thm:incoherence}). The Second Law of Thermodynamics, restated in this framework, is the statement that the partition count is strictly monotone increasing along physical trajectories---a statement that is true by construction, not by statistical accident.
\end{theorem}

\begin{proof}
We summarise the argument. The Time-Count Identity (Theorem~\ref{thm:time-count}) establishes that time is the partition count generated by oscillation, not a parameter independent of the system. The Loop Identity Lemma (Lemma~\ref{lem:loop}) establishes that observable states must carry partition structure beyond microstate coordinates. The Sliding-Endpoint Theorem (Theorem~\ref{thm:sliding}) establishes that the empirical reproducibility of mass spectra is logically equivalent to the irreversibility of partition transitions. The Constitutive Asymmetry Theorem (Theorem~\ref{thm:constitutive}) establishes that physical trajectories are sequences of constitutively one-way operations. The Rewind-as-Forward Principle (Principle~\ref{princ:rewind}) and the Incoherence-of-Reversibility Theorem (Theorem~\ref{thm:incoherence}) establish that no physical operation decrements the partition count; the mathematical automorphism $T$ has no operational realisation. The Source-Analyzer Indeterminacy Theorem (Theorem~\ref{thm:source-indeterminacy}) establishes that the detector recovers only the fixed-point value of the analyzer chain's composed negation operator. The Three-Route Equivalence (Theorem~\ref{thm:three-routes}) establishes that mass, as residue, confinement cost, and negation fixed point, are the same quantity. The Mass-Time-Identity-Count Equivalence (Theorem~\ref{thm:equivalence}) establishes that mass, time, identity, and partition count are unit-conversions of a single quantity.

The composite picture: time is constituted by partition operations; partition operations are constitutively one-way; reversal is not a physical operation; the macroscopic irreversibility we observe is the operational content of what time is, not a statistical regularity layered on top of an underlying time-symmetry. Loschmidt's paradox dissolves because its premise---that microscopic reversibility and macroscopic irreversibility coexist as commensurable claims---is malformed. \qed
\end{proof}

\subsection{Comparison with Existing Resolutions}

\begin{table*}[t]
\centering
\small
\caption{Comparison of resolutions of Loschmidt's paradox.}
\label{tab:comparison}
\begin{tabular}{p{0.18\textwidth}p{0.30\textwidth}p{0.45\textwidth}}
\toprule
\textbf{Approach} & \textbf{Treatment of time-reversal} & \textbf{Limitation} \\
\midrule
Boltzmann statistical \citep{boltzmann1877} & Reversal is possible but improbable, $P \sim e^{-N}$. & Leaves open the in-principle possibility of reversal with sufficient luck or technology. Does not explain reproducibility of single-ion measurements. \\
\addlinespace
Gibbs coarse-graining \citep{gibbs1902} & Reversal is blocked by information loss in coarse-graining. & Makes irreversibility observer-dependent; different coarse-grainings yield different entropies. \\
\addlinespace
Penrose cosmological \citep{penrose1989} & Reversal would require time-symmetric initial conditions; ours are asymmetric. & Explains why entropy is currently increasing but not why local reversal events are absent. \\
\addlinespace
Zurek decoherence \citep{zurek2003} & Reversal is blocked by environmental entanglement. & Applies only to open quantum systems; fails for closed Hamiltonian systems. \\
\addlinespace
Sachikonye partition (\citeyear{sachikonye2025loschmidt}) & Reversal is blocked by non-enumerability of non-actualisations and partition irreversibility. & Operates within the framework that takes time-reversibility as a meaningful concept; resolves the paradox as a within-framework impossibility. \\
\addlinespace
\textbf{Present work} & \textbf{Reversal is a category error: automorphisms of solution sets are not physical operations.} & \textbf{Dissolves the paradox by exposing the malformedness of its premise; supplies empirical anchor through MS reproducibility.} \\
\bottomrule
\end{tabular}
\end{table*}

\subsection{The Empirical Significance of Mass Spectrometry}

A central claim of this paper is that mass spectrometry is not merely an analytical chemistry tool but an empirical demonstration of constitutive asymmetry. We summarise:

\begin{itemize}[leftmargin=1.2em,itemsep=0.2em]
    \item Every reproducible mass spectrum is a verification of the Sliding-Endpoint Theorem.
    \item Every chain of analyzer stages is a sequence of constitutively one-way operations.
    \item Every detection event is the irrecoverable commitment of ion-analyzer correlations to the detector hardware.
    \item Every $m/z$ peak is the operational manifestation of the mass-time-identity-count equivalence.
\end{itemize}

The mass spectrometer is, in this sense, the most direct laboratory witness of the irreversibility of nature. Its routine reproducibility is the operational signature of the Second Law in a form that cannot be argued with on statistical grounds.

\subsection{Implications for the Foundations of Physics}

The resolution proposed has implications beyond the immediate paradox.

\textbf{The status of time.} Time, in this framework, is not a primitive parameter into which dynamics is embedded. It is constituted by partition operations and is unit-equivalent to mass, identity, and partition count. The continuous-time formalism of classical and quantum mechanics is a useful idealisation that obscures the operational structure of time.

\textbf{The status of the Second Law.} The Second Law of Thermodynamics, in this framework, is not a statistical regularity inferred from large-system behaviour. It is the operational content of the partition-counting structure: the partition count cannot decrement because the operations that constitute it are individually one-way. Statistical mechanics provides quantitative refinements, but the qualitative content of the Second Law is structural rather than statistical.

\textbf{The status of measurement.} Measurement, in this framework, is the irreversible commitment of system-instrument correlations to instrumental records. The measured value is the fixed point of a composed negation operator implemented by the instrument. The instrument's specific structure (the operator) is, generically, irrecoverable from the readout (the fixed-point value). This generalises the Source-Analyzer Indeterminacy Theorem from MS to all measurement.

\textbf{The status of identity.} Particle identity, in this framework, is constituted by survival as fixed point of negation through a sequence of partition transitions. It is not preserved by some abstract through-time relation but operationally constituted by the partition chain.

\begin{figure*}[!htbp]
    \centering
    \includegraphics[width=\textwidth]{figures/panel_E6_no_path_backward.png}
    \caption{\textbf{No Path Backward: past states are not future-targets; no operational path returns to $S_0$.}
    (\textbf{A}) Forward trajectory's cumulative partition count over six stages (stage durations drawn uniformly from $[50, 500]$~$\mu$s). The blue line shows $M(\mathrm{stage}\,k) = \sum_{i=1}^{k} f \cdot \tau_i$, monotone-increasing from $M_{S_0} = 0$ at stage $0$ (green dotted) to $M_{S_1} = 18{,}094$~cycles at stage $6$ (red dashed). The trajectory cannot decrement past any prior stage's partition count, by Theorem~\ref{thm:refined-inverse}. The green dotted line marks the unreachable initial-state count.
    (\textbf{B}) Final partition count produced by each of twenty independent attempts to construct a backward operational path from $S_1$ to $S_0$, shown as purple bars. Each attempt consists of a six-step random sequence of inverse operations with durations drawn from $[50, 500]$~$\mu$s. The red dashed line marks $M_{S_1}$ (the forward target); the green dotted line marks $M_{S_0} = 0$ (the supposed reverse target). Every attempt produces a count strictly exceeding $M_{S_1}$, and none reach $M_{S_0}$. The systematic overshoot demonstrates that no operational path can decrement the count.
    (\textbf{C}) Three-dimensional trajectory in (step, $M$, configuration-index) space showing the forward leg (blue circles, $0 \to M_{S_1}$) and the attempted reverse leg (red squares, $M_{S_1} \to 2 M_{S_1}$). The configuration-index axis (vertical) increases through the forward leg (configurations $0 \to 6$) and decreases through the reverse leg (configurations $6 \to 0$), giving the appearance of a back-tracking motion in configuration space; the partition-count axis ($y$), however, increases monotonically through both legs. The green star at the origin marks $S_0$, the unreachable target. The plot visually unifies the configurational return (which exists) with the partition-count irreversibility (which prevents true return).
    (\textbf{D}) Final partition count for fifty random forward trajectories (each with three to nine stages, durations drawn from $[10, 1000]$~$\mu$s), plotted on the trajectory index axis with plasma colourmap encoding the count value. All fifty trajectories terminate with $M > 0$, and none admit a backward operational path that would decrement the count to $M_{S_0} = 0$ (green dotted). The complete absence of backward paths across the random sample is the structural validation of the No Path Backward Theorem.}
    \label{fig:panel-E6}
\end{figure*}

%==============================================================================
\section{Numerical Validation}
\label{sec:validation}
%==============================================================================

We now present a numerical validation of the principal theorems of the paper. Seven experiments were conducted, each testing a distinct theorem or principle, with a total of 162 sub-tests across approximately five orders of magnitude of physical parameters. All 162 sub-tests passed; all seven experiments achieved full passing status. The validation suite, raw JSON results, and figure-generation scripts are provided in the supplementary materials.

\subsection{Experiment 1: Time-Count Identity}
\label{sec:val-E1}

The Time-Count Identity (Theorem~\ref{thm:time-count}) was tested over $11$ time scales spanning $10$ orders of magnitude (1~ns to 10~s) and $6$ oscillator frequencies spanning $12$ orders of magnitude (1~kHz to 1~PHz), giving $66$ test conditions. For each $(f, t)$ pair, the cycle count $M = \lfloor f t \rfloor$ was computed, and the recovered time $t_{\mathrm{rec}} = M/f$ was compared with the input time.

\textbf{Result:} For all conditions in which $f t \geq 1$ (i.e., at least one full cycle accumulated), the relative reconstruction error $|t_{\mathrm{rec}} - t|/t$ was bounded by $10^{-15}$, the floating-point precision of the computation. The conditions for which $f t < 1$ correspond to sub-cycle observation windows, where the count is unable to resolve fractional cycles---a regime in which the Time-Count Identity is operationally meaningful only at the resolution $1/f$. The mean relative error in the resolvable regime was $3.2 \times 10^{-17}$. All $66/66$ tests passed.



\subsection{Experiment 2: Sliding-Endpoint Theorem}
\label{sec:val-E2}

The Sliding-Endpoint Theorem (Theorem~\ref{thm:sliding}) was tested by simulating $100$ ions distributed across $m/z = 100$--$1500$, each measured $5$ times under identical conditions to assess reproducibility, and then varying the stop time across $11$ values to confirm the count's monotonic dependence on $t_C$. A counterfactual test computed the hypothetical mass-shift that would result from deleting partition transitions, contrasting it against the empirically observed zero shift.

\textbf{Result:} The replicate spread across all reproducibility tests was bounded by $1.14 \times 10^{-16}$ (floating-point precision). The cycle count was strictly monotone-increasing across the $11$ stop times. The maximum hypothetical mass-shift from a $\Delta M = 1000$ deletion would be $\sim 50$~Da---a shift that is empirically forbidden by the irreversibility of partitions. All $25/25$ tests passed.



\subsection{Experiment 3: Rewind-as-Forward Principle}
\label{sec:val-E3}

The Rewind-as-Forward Principle (Principle~\ref{princ:rewind}) and the Refined Inverse Theorem (Theorem~\ref{thm:refined-inverse}) were tested by simulating a 5-stage forward analyzer chain ($M_{\mathrm{forward}} = 4{,}100$ cycles) and then applying a putative ``reversal'' operation. The total partition count after the round trip was compared with the forward-only count. A second test verified that any inverse operation increments the partition count by at least one cycle, across $20$ randomly-generated inverse operations. A third test repeated the round-trip simulation $30$ times with random stage compositions.

\textbf{Result:} After the reversal, the partition count was $M_{\mathrm{after}} = 8{,}200$ cycles---exactly twice the forward count, as predicted by Theorem~\ref{thm:refined-inverse}. All $20$ inverse operations incremented the count by $\geq 1$~cycle, with the smallest increment being $1.0$~cycles ($100$~ns at $10$~MHz). Across $30$ random round-trip trials, $30/30$ produced strictly larger final counts than their forward-only counterparts. All $22/22$ tests passed.



\subsection{Experiment 4: Source-Analyzer Indeterminacy}
\label{sec:val-E4}

The Source-Analyzer Indeterminacy Theorem (Theorem~\ref{thm:source-indeterminacy}) was tested by constructing $30$ distinct analyzer chains, each containing between $2$ and $7$ stages with random stage durations, all configured to produce the same total partition count and the same $m/z$ readout ($m/z = 500$). The pairwise distinctness of the chains was assessed, and the source information committed to ion-analyzer correlations was quantified using Shannon entropy of the stage-duration distribution.

\textbf{Result:} All $30$ chains produced exactly the same readout to within floating-point precision. Of $\binom{30}{2} = 435$ chain-pairs, $435/435$ were operationally distinct (different stage counts or different stage-duration distributions). The mean Shannon source-entropy per chain was $3.544$~bits, corresponding to $\sim 3.4 \times 10^{-23}$~J/K of categorical entropy committed to correlations and unrecoverable from the detector readout. All $3/3$ tests passed.



\subsection{Experiment 5: Mass-Time-Identity-Count Equivalence}
\label{sec:val-E5}

The Mass-Time-Identity-Count Equivalence (Theorem~\ref{thm:equivalence}) was tested on a panel of $10$ representative ions spanning $4$ orders of magnitude in mass ($1$~Da, proton, to $14{,}305$~Da, lysozyme-10+). For each ion, the rest mass was converted to rest frequency via $\omega_0 = mc^2/\hbar$, the rest energy was independently computed, and the equality $E_{\mathrm{rest}} = \hbar \omega_0$ was verified. The Time-Count Identity $t = M/f$ was independently verified for each ion at a $1$~ms observation window.

\textbf{Result:} Across the ion panel, the mean relative error in the energy-frequency identity $E = \hbar\omega_0$ was $3.2 \times 10^{-17}$, dominated entirely by floating-point round-off. The mean relative error in $t = M/f$ was below $10^{-12}$. The four quantities (mass, frequency, period, count) were verified to satisfy the unit-conversion relations exactly. All $31/31$ tests passed.



\subsection{Experiment 6: No Path Backward}
\label{sec:val-E6}

The No Path Backward Theorem (Theorem~\ref{thm:no-path-backward}) was tested by simulating a $6$-stage forward trajectory (cumulative count $M_{S_1} = 18{,}094$ cycles), and then attempting $20$ distinct ``backward path'' constructions, each consisting of a random sequence of inverse operations that purport to return the system from $S_1$ to $S_0$. The candidate state's partition count was compared with the true $M_{S_0} = 0$. A separate test ran $50$ random forward trajectories and counted how many admitted a valid backward operational path.

\textbf{Result:} Of $20$ backward-path attempts, $20/20$ produced final partition counts strictly exceeding $M_{S_1}$, never reaching $M_{S_0}$. The candidate ``past state'' produced by the round trip had $M_{\mathrm{candidate}} = 36{,}189$ cycles, vastly exceeding the true $M_{S_0} = 0$. The becoming-history of the candidate state had length $12$ (forward + reverse), while the true $S_0$ had length $0$. Of $50$ random forward trajectories, $50/50$ admitted no valid backward path. All $4/4$ tests passed.



\subsection{Experiment 7: Oscillation as Forward Re-Placement}
\label{sec:val-E7}

The Oscillation as Forward Re-Placement Principle (Principle~\ref{princ:oscillation}) and the Cycle Distinguishability Theorem (Theorem~\ref{thm:cycle-distinguishability}) were tested by simulating a 1~MHz harmonic oscillator $x(t) = A\cos(\omega t)$ and sampling its position at $101$ cycle peaks. The configurational identity of all peak-positions was verified, and the partition count at each peak was independently computed via $M = ft$. A long-run test extended the simulation to $10{,}000$ cycles to confirm that the count does not reset.

\textbf{Result:} At all $101$ peak times, the position was identically $A$ (configurational return) to within floating-point precision; however, the partition counts at these times were all distinct and strictly monotone-increasing. The count increment per cycle was exactly $f \cdot T = 10$ (consistent with $f = 10$~MHz reference oscillator and $T = 1~\mu$s ion-oscillator period). After $10{,}000$ cycles, the partition count was $M = 99{,}990$, monotone throughout---five orders of magnitude larger than the naive ``return'' hypothesis would predict. The cycle distinguishability $\Delta M = k$ was verified exactly for $k \in \{1, 2, 5, 10, 100, 1000\}$. All $11/11$ tests passed.

\begin{figure*}[!htbp]
    \centering
    \includegraphics[width=\textwidth]{figures/panel_E7_oscillation_forward.png}
    \caption{\textbf{Oscillation as Forward Re-Placement: configurational return without operational return, verified across $10{,}000$ cycles.}
    (\textbf{A}) Position $x(t) = A \cos(\omega t)$ of a $1$~MHz harmonic oscillator over eight cycles (blue continuous curve), with red points marking the cycle peaks at times $t_n = n \cdot T = n \cdot 1$~$\mu$s for $n = 0, 1, \ldots, 8$. At every peak, the position is exactly $A$: the configurations are identical to floating-point precision. This is the configurational return that, on the naive view, would correspond to ``the oscillator going back to where it was.''
    (\textbf{B}) Partition count $M(t) = f \cdot t$ of the same oscillator over the same eight cycles (green continuous curve), with red points at the cycle peaks. Despite the configurational return shown in (A), the partition count at each peak is strictly greater than at the previous peak, increasing linearly by $\Delta M = f \cdot T = 10$ cycles per period (with $f = 10$~MHz and $T = 1$~$\mu$s). The strict monotonicity is the operational content of forward re-placement: each peak is a new instance whose identity is constituted by its larger partition count.
    (\textbf{C}) Three-dimensional trajectory in $(t, x, M)$ space over eight cycles, with blue continuous curve tracing the oscillator's joint motion and viridis-coloured points at the cycle peaks. The trajectory forms an ascending helix: the position $x$ cycles between $\pm A$ (radial axis), but the partition count $M$ ascends monotonically (vertical axis). The ascending-helix geometry is the visual signature of oscillation as forward re-placement: configurational periodicity coexists with operational forward-only motion. The viewing angle (elevation $20^\circ$, azimuth $-60^\circ$) maximises visibility of the helical structure.
    (\textbf{D}) Long-run partition count over $10{,}000$ cycles, plotted on a symmetric-logarithmic vertical axis. The orange continuous curve traces $M(\mathrm{cycle\,index})$ from $M = 0$ at index $0$ to $M = 99{,}990$ at index $10{,}000$, monotonically increasing throughout. The dashed red line at $M = 1$ marks the maximum count predicted by the naive ``return'' hypothesis (the oscillator's partition count would oscillate between $0$ and $1$ if cycles were genuine returns to past states). The empirical count exceeds this naive bound by five orders of magnitude, providing direct empirical falsification of the return hypothesis. The result confirms Principle~\ref{princ:oscillation}: oscillation is not return but forward re-placement.}
    \label{fig:panel-E7}
\end{figure*}


\subsection{Aggregate Validation Results}
\label{sec:val-aggregate}

\begin{table}[H]
\centering
\caption{Summary of validation experiments. All $162$ sub-tests passed across $7$ experiments testing the principal theorems of the paper.}
\label{tab:validation-summary}
\small
\begin{tabular}{cp{0.40\columnwidth}rr}
\toprule
\textbf{ID} & \textbf{Theorem} & \textbf{Tests} & \textbf{Status} \\
\midrule
E1 & Time-Count Identity (Thm~\ref{thm:time-count}) & 66/66 & PASS \\
E2 & Sliding-Endpoint (Thm~\ref{thm:sliding}) & 25/25 & PASS \\
E3 & Rewind-as-Forward (Thm~\ref{thm:refined-inverse}) & 22/22 & PASS \\
E4 & Source Indeterminacy (Thm~\ref{thm:source-indeterminacy}) & 3/3 & PASS \\
E5 & Mass-Time-Count-Identity (Thm~\ref{thm:equivalence}) & 31/31 & PASS \\
E6 & No Path Backward (Thm~\ref{thm:no-path-backward}) & 4/4 & PASS \\
E7 & Oscillation as Forward (Thm~\ref{thm:cycle-distinguishability}) & 11/11 & PASS \\
\midrule
\textbf{Total} & & \textbf{162/162} & \textbf{7/7} \\
\bottomrule
\end{tabular}
\end{table}

The validation suite provides quantitative empirical support for every principal theorem. Identities involving fundamental physical constants ($t = M/f$, $E = \hbar\omega_0$) hold to floating-point precision across many orders of magnitude. Operational claims (no decrement of $M$, no backward path, monotone-increasing $M$ through cycles) are verified across hundreds of randomised trials. Information-theoretic claims (chain entropy, source information committed to correlations) are quantified to bit-level precision. We conclude that the framework is internally consistent, numerically validated, and ready for physical-data tests on instrument records.

%==============================================================================
\section{Discussion}
\label{sec:discussion}
%==============================================================================

\subsection{Objections and Responses}

We address several anticipated objections.

\textbf{Objection 1.} \emph{The argument depends on specific properties of mass spectrometry. Why should a result derived from a particular instrument apply to physics generally?}

\emph{Response.} The argument uses MS as an empirical witness, not as a special case. The Time-Count Identity, Loop Identity Lemma, Constitutive Asymmetry Theorem, and Incoherence Theorem are framework-level results that apply to any bounded oscillatory system. MS is the most direct empirical instantiation, but the structural content is general.

\textbf{Objection 2.} \emph{Reversibility is well-defined in classical mechanics: $\dot{q}_i = \partial H/\partial p_i$ admits time-reversed solutions. How can it be a category error?}

\emph{Response.} The classical-mechanics reversibility is a property of the solution set under the mathematical operation $T: x(t) \to x(-t)$. It is well-defined as an automorphism of the solution set. The category error is the additional claim that this automorphism corresponds to a physical operation that could be enacted in the laboratory. There is no such enactment; no laboratory procedure decrements the partition count.

\textbf{Objection 3.} \emph{Quantum mechanics has unitary evolution, which is reversible by construction. How does the framework apply?}

\emph{Response.} Quantum unitary evolution describes the action of the Schrödinger operator on state vectors. Like classical Hamiltonian flow, it is an automorphism of the solution set. The application of the inverse operator is itself a quantum operation that takes physical time and generates partition transitions; it does not decrement the global partition count. Quantum measurement, moreover, is an explicitly irreversible operation \citep{vonneumann1932}, and the partition framework treats it as a partition-creating event \citep{sachikonye2025loschmidt}.

\textbf{Objection 4.} \emph{The argument seems to rule out ``time-reversed simulations'' on classical computers, but these clearly run.}

\emph{Response.} A time-reversed simulation on a classical computer is an entirely forward-time process. The simulation's clock advances forward; the computer's CPU dissipates heat forward; the simulator's wristwatch ticks forward. The simulation's content---a representation of states traversed in reverse order---is not a reversal of physical time. By the Rewind-as-Forward Principle, it is a forward operation whose output happens to be a reverse-ordered list of states.

\textbf{Objection 5.} \emph{The framework appears to make time-reversibility unfalsifiable. Doesn't this make it unscientific?}

\emph{Response.} The framework does not make time-reversibility unfalsifiable; it makes it ill-formed. An ill-formed claim is not a claim at all, just as ``the unmarried bachelor is married'' is not a falsifiable claim about the world. The empirical content of the framework lies elsewhere: in the predictions about mass spectrometry, partition counting, and trans-Planckian resolution \citep{sachikonye2025temporal}, all of which are empirically falsifiable.

\subsection{Connections to Other Foundational Issues}

The resolution proposed connects to several other foundational issues in physics.

\textbf{The measurement problem in quantum mechanics.} The conflict between unitary evolution and measurement collapse is a special case of the conflict between automorphisms-of-solution-sets and physical operations. Unitary evolution is the former; measurement is the latter. The conflict dissolves once we recognise that they live at different categorical levels.

\textbf{Black hole information.} Hawking's information paradox \citep{hawking1976} arises from assuming that information is preserved by unitary evolution but apparently lost in black hole evaporation. The partition framework suggests that the information is preserved as partition residue accumulated on the horizon \citep{sachikonye2025loschmidt}. The conflict between ``unitarity'' (an automorphism property) and ``information loss'' (an operational claim) is structurally similar to Loschmidt's paradox.

\textbf{The cosmological arrow of time.} The framework provides a structural rather than initial-conditions account of the cosmological arrow of time. Time flows forward because partition operations are individually one-way, not because the universe began in a special state.

\textbf{Computability and physics.} The Constitutive Asymmetry Theorem suggests that physical processes are not equivalent to abstract dynamical systems. They are sequences of operations, each with mechanism-specific realisations. This suggests connections to algorithmic and computational accounts of physics.

\subsection{Predictions and Empirical Tests}

The framework makes several empirically testable predictions.

\textbf{Prediction 1: MS reproducibility scales with instrument count integrity.} The reproducibility of mass spectra should depend on the integrity of the partition counting. Instruments that maintain coherent partition counting through the chain should produce more reproducible spectra than those that do not. This is testable by comparing spectra across instruments with different counting precision.

\textbf{Prediction 2: Categorical entropy generation per analyzer stage is bounded below.} By Equation~\eqref{eq:chain-entropy}, each analyzer stage generates at least $\kB \ln 2$ of categorical entropy. This sets a fundamental limit on the entropy efficiency of mass spectrometry that is independent of mechanical or thermal sources.

\textbf{Prediction 3: Source indeterminacy is structurally bounded.} Different analyzer chains producing identical readouts should be empirically distinguishable only through correlations external to the readout itself (e.g., differential heat dissipation, electronic noise floor). This predicts a measurable amount of source-information committed to ion-analyzer correlations.

\textbf{Prediction 4: Trans-Planckian temporal resolution.} By Theorem 6.1 of \citep{sachikonye2025temporal}, the partition resolution scales as $\delta t = T/\Mcount$, with no quantum-mechanical lower bound. This predicts that mass spectrometers achieve effective temporal resolution far beyond the oscillator period, which is empirically observed.

%==============================================================================
\section{Conclusion}
\label{sec:conclusion}
%==============================================================================

We have resolved Loschmidt's paradox by exposing the malformedness of its premise. The paradox is generated by treating the mathematical automorphism $T: x(t) \to x(-t)$ as if it were a physical operation. It is not. Physical operations are constitutively one-way; their composite trajectories have no reverse trajectories; any operation purporting to enact a reversal is itself a forward operation that increments the partition count. Time-reversibility is therefore not a property dynamics either has or lacks but a category error generated by mistaking automorphisms of solution sets for physical operations on systems.

The empirical anchor of this resolution is mass spectrometry. The reproducibility of mass spectra is logically equivalent to the irreversibility of partition transitions. Every reproducible spectrum is a written record of constitutive asymmetry. The mass spectrometer, far from being a specialised analytical chemistry tool, is the most direct laboratory witness of the operational asymmetry of physical reality.

The framework collapses four traditionally distinct quantities---mass, time, identity, and partition count---into a single unit-equivalent observable. Mass is what we read off the instrument. Time is what the instrument tells. Identity is what survives the analyzer chain. Partition count is what increments through the chain. The four are coordinatisations of a single physical quantity. The mass spectrum is the elapsed proper time of its constituent particles, expressed in mass units. The arrow of time is made visible at every detection event.

The Second Law of Thermodynamics, in this view, is not a statistical regularity. It is the operational signature of the constitutive directionality of physical operations. It does not need to be inferred from large-system behaviour; it is built into what time is.

The maize seed in your hand was not arrived at by reversing time. It was constituted by a sequence of one-way operations whose composite produced the seed. The ion at the detector was not arrived at by reversing time. It was constituted by a sequence of one-way operations whose composite produced the spectrum. In both cases, the very concept of reaching the end-state by reversing the process is incoherent. There is nothing to reverse, and no operation that would reverse it.

Loschmidt asked: if the laws are reversible, why is the world irreversible? The answer is: the laws were never reversible in the operational sense. The reversibility was a feature of the equations as automorphisms of their solution sets. The physical processes the equations approximate are constitutively one-way. The world is irreversible because its operations are.

%==============================================================================
\section*{Acknowledgements}
%==============================================================================

The author thanks the development of the partition counting framework \citep{sachikonye2025loschmidt,sachikonye2025partition,sachikonye2025counting,sachikonye2025temporal} which provides the technical apparatus deployed in this paper. The empirical validation discussed here builds on the experimental work reported in those papers.

%==============================================================================
\bibliographystyle{plainnat}
\bibliography{references}
%==============================================================================

\end{document}